\par\noindent\textbf{(1)} 一个往返相位为 \(\Phi(\omega)=2n(\omega)\omega L/c+\phi_1(\omega)+\phi_2(\omega)\)。谐振条件 \(\Phi(\omega_m)=2\pi m\)。若色散弱，可将它近似为整数光学长度条件；强色散镜的反射相位绝不能随意置零。
\par\noindent\textbf{(2)} 对相邻纵模条件差分 \(2\pi\approx\Phi'(\omega_m)\Delta\omega\)，所以 \(\Delta\omega_{\rm FSR}\approx2\pi/\tau_g\)，其中 \(\tau_g=\Phi'(\omega)=2L[n+\omega n']/c+\phi_1'+\phi_2'\) 是总往返群延迟。由此可用色散镜调节模间距。
\par\noindent\textbf{(3)} 设每往返复场幅衰减模为 \(0<r<1\)，响应几何级数分母 \(1-re^{i\Phi}\)。共振近旁强度分母 \((1-r)^2+r[\tau_g(\omega-\omega_m)]^2\)，故角频率全宽 \(\delta\omega_{\rm FWHM}\approx2(1-r)/(\sqrt r|\tau_g|)\)，在 \(r\to1\) 时线宽趋窄。
\par\noindent\textbf{(4)} 镜面局部相位导数 \(d\phi_i/d\omega\) 可以为负，但总输出脉冲群延迟 \(\tau_g=d\Phi/d\omega\) 不能仅靠一个频点斜率推断信息超光速；有限带宽信号仍由因果响应约束。强吸收、明显群延迟色散与有限带宽时，线宽的Lorentz近似失效，应使用完整复反射系数 \(r_i(\omega)\) 与因果解析结构。
